\documentclass[border=14pt]{standalone}
\usepackage{ctex}
\usepackage{tikz}
\begin{document}
\begin{tikzpicture}[x=1cm, y=1cm, font=\normalsize, line width=0.8pt]

% fig-ge-03 R / L / C 相量伏安（中文；公式写图注不进图内）
% 布局：每栏 5.2 宽，原点在 (0.6,1.6)；标签放箭头末端外侧；角度标注在弧正上方/下方

\def\cw{5.2}
\def\gap{1.6}

% 栏宏：#1 x 位移 #2 栏名 #3 说明 #4 主体
\newcommand{\sub}[4]{%
  \begin{scope}[shift={(#1,0)}]
    \node[anchor=south] at ({\cw/2},4.3) {#2};
    \draw[->] (0.6,1.6) -- ({\cw-0.5},1.6) node[anchor=west, font=\small] {Re};
    \draw[->] (0.6,1.6) -- (0.6,3.9) node[anchor=south, font=\small] {Im};
    #4
    \node[anchor=north, font=\small] at ({\cw/2},-1.30) {#3};
  \end{scope}}

% ── R：同相 ──
\sub{0}{R}{U 与 I 同相}{
  \draw[->, line width=1.4pt] (0.6,1.6) -- ++(2.8,1.3);
  \node[anchor=south west, font=\small] at (0.6+2.8,1.6+1.3) {U};
  \draw[->, line width=0.9pt, dashed] (0.6,1.6) -- ++(2.1,0.98);
  \node[anchor=north west, font=\small] at (0.6+2.1,1.6+0.98) {I};
}

% ── L：I 滞后 U 90 deg ──
\sub{{\cw+\gap}}{L}{}{
  \draw[->, line width=1.4pt] (0.6,1.6) -- ++(2.8,1.3);
  \node[anchor=south west, font=\small] at (0.6+2.8,1.6+1.3) {U};
  % I：把 U 方向转 -90 → (1.3, -2.8)*0.45
  \draw[->, line width=0.9pt, dashed] (0.6,1.6) -- ++(1.25,-2.0);
  \node[anchor=north east, font=\small] at (0.6+1.10,1.6-1.95) {I};
  % 角度弧：从 U 方向到 I 方向
  \draw[->, line width=0.5pt] (1.9,2.2) arc (24.6:-58:0.9);
  \node[anchor=north east, font=\small] at ({\cw/2},-1.30) {I 滞后 U};
  \node[anchor=north west, font=\small] at ({\cw/2+0.06},-1.30) {90°};

}

% ── C：I 超前 U 90 deg ──
\sub{{2*(\cw+\gap)}}{C}{}{
  \draw[->, line width=1.4pt] (0.6,1.6) -- ++(2.8,1.3);
  \node[anchor=south west, font=\small] at (0.6+2.8,1.6+1.3) {U};
  % I：把 U 方向转 +90 → (-1.3, 2.8)*0.45
  \draw[->, line width=0.9pt, dashed] (0.6,1.6) -- ++(-1.25,2.0);
  \node[anchor=south east, font=\small] at (-0.65,3.6) {I};
  \draw[->, line width=0.5pt] (1.9,2.2) arc (24.6:114.6:0.9);
  \node[anchor=north east, font=\small] at ({\cw/2},-1.30) {I 超前 U};
  \node[anchor=north west, font=\small] at ({\cw/2+0.06},-1.30) {90°};

}

\end{tikzpicture}
\end{document}